\chapter{Vacuum Energy, Inflation, and Primordial Fluctuations}

\section{Vacuum Energy and de Sitter Space}

In flat space the Friedmann equation is
\begin{equation}
\left(\frac{\dot A}{A}\right)^2=\frac{8\pi G}{3}\rho.
\end{equation}
For matter, \(\rho\) decreases like \(A^{-3}\). For radiation, \(\rho\) decreases like \(A^{-4}\). For vacuum energy it does not dilute:
\begin{equation}
\rho=\rho_0=\text{constant}.
\end{equation}
If the universe contains nothing but this vacuum energy, the right-hand side is a constant. Writing that constant as \(H^2\),
\begin{equation}
H^2=\frac{8\pi G}{3}\rho_0,
\qquad
H=\sqrt{\frac{8\pi G}{3}\rho_0}.
\end{equation}
Then the Hubble constant is truly constant:
\begin{equation}
\frac{\dot A}{A}=H.
\end{equation}
So
\begin{equation}
\dot A=HA,
\end{equation}
and the solution is exponential:
\begin{equation}
A(t)=C e^{Ht}.
\end{equation}
The distance between distant unbound galaxies grows exponentially. This is de Sitter space.

In these coordinates the geometry is
\begin{equation}
d\tau^2=dt^2-e^{2Ht}\left(dx^2+dy^2+dz^2\right).
\end{equation}
The vacuum energy here is positive. In flat space it has to be, because the left-hand side of the Friedmann equation is nonnegative. A negative vacuum energy can appear only when curvature terms are included, and that is not the case being followed here.

Einstein introduced the cosmological constant for a different purpose: to balance ordinary matter and hold the universe static.

\section{The Present Accelerating Universe}

De Sitter space is also the late-time tendency of the present universe. Earlier the scale factor followed the matter-dominated behavior
\begin{equation}
A(t)\propto t^{2/3}.
\end{equation}
A few billion years ago the curve bent upward and began to increase more rapidly. It is not yet a perfect exponential, because other forms of energy are still present, but the trend is toward one.

Acceleration here means
\begin{equation}
\ddot A>0.
\end{equation}
This is an observed feature of the late-time expansion.

Now consider the horizon. Hubble's law says
\begin{equation}
v=HD.
\end{equation}
At the distance where the recession speed becomes the speed of light,
\begin{equation}
D=\frac{c}{H}.
\end{equation}
Equivalently, one can rewrite this as \(HD=c\).\editorialnote{This is an editorial algebraic restatement of the spoken relation \(D=c/H\). The upper line on the right board in Figure~\ref{fig:lecture07-horizon} is only partly legible.}

\lecturetimestamp{00:30:37}
\begin{figure}[ht]
\centering
\includegraphics[width=0.62\textwidth]{lecture_07_figure_02.png}
\caption{Cosmic horizon relation \(D=c/H\) in de Sitter expansion.}
\label{fig:lecture07-horizon}
\end{figure}

This distance is of order ten to fifteen billion light years, numerically close to the age of the universe once the latter is converted to a distance with \(c\).

\begin{classroomqa}
\lecturetimestamp{00:14:10}
\audiencequestion{Isn't the near equality between the horizon scale and the age of the universe just what happens if the tangent line happens to meet the origin?}
\lecturerresponse{That would be exactly true for \(A(t)\propto t\). It is not true in general. In de Sitter space \(H\) stays fixed while the age keeps increasing, so the present near equality is not a mathematical identity.}
\end{classroomqa}

\section{What Vacuum Energy Means}

Vacuum energy is described here in the language of ordinary quantum mechanics. For a harmonic oscillator the energy is schematically
\begin{equation}
E\sim p^2+x^2.
\end{equation}
Classically the ground state is the equilibrium point,
\begin{equation}
x=0,\qquad p=0,\qquad E=0.
\end{equation}
Quantum mechanics does not allow both \(x\) and \(p\) to be exactly fixed at zero. The uncertainty principle forces fluctuations, so the ground-state energy is positive. For the oscillator one gets
\begin{equation}
E_0=\frac{1}{2}\hbar\omega.
\end{equation}

The same point is then carried over to fields. The universe is full of harmonic oscillators: the oscillations of the electromagnetic field, and likewise of every other field. Each field has zero-point energy. For the electromagnetic field the energy density is schematically
\begin{equation}
u_{\mathrm{EM}}\sim E^2+B^2.
\end{equation}
Even with no real photons present there remain irreducible zero-point oscillations. That is one explanation of vacuum energy. The only property needed here is that this energy is present and does not dilute as the universe expands.

\begin{classroomqa}
\lecturetimestamp{00:16:15}
\audiencequestion{What is this vacuum energy, and why would anything remain if matter and radiation were removed?}
\lecturerresponse{Because the vacuum is the ground state of the quantum fields, and a quantum ground state need not have zero energy. Harmonic oscillators have positive ground-state energy, and fields behave as collections of oscillators. So even with no real particles present there remains zero-point energy.}
\end{classroomqa}

This is not the three-degree background radiation. The microwave background is ordinary radiation. Its energy density decreases because the photons dilute and each photon redshifts. Vacuum energy is the residual zero-point energy that remains after the real photons have been removed.

\section{What \(H^{-1}\) Means}

Now return to the relation between the inverse Hubble constant and the age of the universe. If
\begin{equation}
A(t)\propto t,
\end{equation}
then
\begin{equation}
H=\frac{\dot A}{A}=\frac{1}{t}.
\end{equation}
More generally, if
\begin{equation}
A(t)\propto t^p,
\end{equation}
then
\begin{equation}
H=\frac{p}{t}.
\end{equation}
For the matter-dominated case \(p=2/3\),
\begin{equation}
H=\frac{2}{3t}.
\end{equation}
So in power-law cosmologies \(H^{-1}\) is of the same order as the age, but not exactly the age.

De Sitter space is different. There
\begin{equation}
A(t)=C e^{\gamma t},
\end{equation}
and therefore
\begin{equation}
H=\frac{\dot A}{A}=\gamma,
\end{equation}
a true constant. Here \(H^{-1}\) is not the age of the universe. It is the time constant for the exponential growth: the time in which the scale factor is multiplied by \(e\). The spoken discussion repeatedly says ``double'' and then corrects it; the precise statement is multiplication by \(e\). Today that time constant is of order ten billion years.

\begin{classroomqa}
\lecturetimestamp{00:23:46}
\audiencequestion{Could you go over the point about the age of the universe and the inverse Hubble constant?}
\lecturerresponse{For power laws one gets \(H=p/t\), so \(H^{-1}\) tracks the age up to a factor of order one. But in de Sitter space \(H\) is constant while the age keeps growing, so \(H^{-1}\) is only the exponential time constant, not the age.}
\end{classroomqa}

\section{Energy and the Cosmic Horizon}

If the vacuum-energy density stays constant while the universe expands, it is natural to ask whether the total vacuum energy is increasing. The answer given here is limited. In general relativity the energy bookkeeping includes a gravitational component, so the notion of total energy is subtler than it is in elementary mechanics. The practical point is the one tied to the horizon.

In exact de Sitter space the cosmic horizon has fixed size,
\begin{equation}
D=\frac{c}{H}
=\frac{c}{\sqrt{8\pi G\rho_0/3}}.
\end{equation}
The volume inside that horizon is therefore fixed. The vacuum energy within it is
\begin{equation}
E_{\mathrm{hor}}=\frac{4\pi}{3}D^3\rho_0.
\end{equation}
\editorialnote{This condenses the spoken calculation at 00:31:06--00:32:23: take the horizon size, form the spherical volume, and multiply by \(\rho_0\).}
Since both \(D\) and \(\rho_0\) are constant, this quantity does not change with time. For the portion of the universe that can ever be seen by a given observer, the accessible vacuum energy stays fixed in exact de Sitter space.

Things can also pass through the cosmic horizon in much the same sense that they pass through a black-hole horizon. Once outside, they are out of causal contact with us.

\begin{classroomqa}
\lecturetimestamp{00:29:27}
\audiencequestion{If the vacuum-energy density stays constant while the universe expands, where does the extra energy come from?}
\lecturerresponse{The lecture answers in general-relativistic language, not Newtonian box language. In exact de Sitter space the horizon size stays fixed, so the vacuum energy inside the region that can ever be seen stays fixed as well.}
\end{classroomqa}

\section{Why Inflation Was Introduced}

The same de Sitter mechanism is then pushed into the remote past, but with a much shorter time constant, of order
\begin{equation}
H^{-1}\sim 10^{-32}\,\mathrm{s}.
\end{equation}
That is the idea of inflation: before anything directly observable, the universe underwent a period of exponential expansion.

The basic motivation is the flatness and smoothness of the universe. The spatial slices appear very flat. On large scales the universe is also very homogeneous. More strikingly, the surface of last scattering was already smooth, so the smoothing could not have happened only during the later visible history.

A deflated balloon is wrinkled. Blow it up enough and any visible patch becomes smooth. The raisin or prune analogy makes the same point: a wrinkled object, stretched by an enormous factor, looks smooth on any visible patch. Inflation is supposed to have done exactly that.

This also answers the horizon problem in the form used here. Widely separated regions of the last-scattering surface are both visible to us, but light did not have time to travel from one to the other during the later visible history. Their nearly uniform density and temperature therefore suggest a prehistory in which one smaller region was stretched enormously.

Inflation has a mixed virtue. It wipes out sensitivity to the initial conditions, and for the same reason it wipes out much of the information about how the universe started.

The flatness question reappears here in sharper form. Observationally the curvature parameter \(k\) looks like zero, at the level of about one percent or better. Geometry looks Euclidean over the visible region. Inflation explains this by making the radius of curvature much larger than the region we can see.

\begin{classroomqa}
\lecturetimestamp{00:40:05}
\audiencequestion{Are we assuming equilibrium before inflation?}
\lecturerresponse{No. The point of inflation is almost the opposite. Whatever structure was there beforehand gets stretched so much that the visible region loses memory of the earlier details.}
\end{classroomqa}

\section{A Scalar Field and Its Potential}

To explain how the effective vacuum energy could have been much larger in the early universe than it is now, introduce a scalar field. Unlike the electric and magnetic fields, which are vectors, a scalar field has no direction. It simply has a value at each point:
\begin{equation}
\phi=\phi(x,y,z,t).
\end{equation}
This field is hypothetical. It is not identified with any known field in nature, and its quanta are not identified with any known particles.

The part of its energy that matters here is a potential energy,
\begin{equation}
V(\phi).
\end{equation}
By ``potential energy'' one means energy associated with the value of the field itself, not with the time derivative of the field. Different values of \(\phi\) can have different energies.

\lecturetimestamp{00:47:08}
\begin{figure}[ht]
\centering
\includegraphics[width=0.68\textwidth]{lecture_07_figure_03.png}
\caption{Qualitative scalar potential \(V(\phi)\) showing that different field values carry different potential energies.}
\label{fig:lecture07-generic-potential}
\end{figure}

\begin{classroomqa}
\lecturetimestamp{00:46:34}
\audiencequestion{What is the horizontal axis in the \(V(\phi)\) drawing?}
\lecturerresponse{The horizontal axis is the field value \(\phi\), and the vertical axis is the potential energy \(V\).}
\end{classroomqa}

If the field is constant in space and time, or nearly constant over the region of interest, then \(V(\phi)\) behaves exactly like vacuum energy. It does not dilute. So the earlier de Sitter formulas are reused with \(\rho_0\) replaced by \(V(\phi)\). In particular,
\begin{equation}
H=\sqrt{\frac{8\pi G}{3}V(\phi)}.
\end{equation}

If \(\phi\) varies from place to place, different regions inflate with different time constants. But if the stretching is large enough, any visible patch becomes nearly homogeneous.

\begin{classroomqa}
\lecturetimestamp{01:05:11}
\audiencequestion{Is \(\phi\) acting like a time parameter here?}
\lecturerresponse{In this discussion it does, in the limited sense that while the field rolls down the potential it is a monotonic function of time.}
\end{classroomqa}

\section{A Plateau, a Cliff, and Reheating}

The potential needed for this story becomes very flat, though not perfectly flat, then reaches the edge of a cliff, drops sharply, does not quite go to zero, and rises again.

\lecturetimestamp{01:02:30}
\begin{figure}[ht]
\centering
\includegraphics[width=0.68\textwidth]{lecture_07_figure_04.png}
\caption{Inflationary potential with a high plateau, a steep drop, and a lower minimum.}
\label{fig:lecture07-inflationary-potential}
\end{figure}

This shape is highly contrived, but it is introduced because the cosmological story seems to require something like it.

The field begins high on the plateau. Because the slope is small and the effective viscosity associated with the expansion is large, \(\phi\) changes only very slowly. A rock falling slowly through honey is closer to the intended picture than an object sliding on a table. The faster the universe expands, the larger this viscosity becomes.

\begin{classroomqa}
\lecturetimestamp{00:59:07}
\audiencequestion{What is causing \(\phi\) to move?}
\lecturerresponse{The slope of the potential is pushing it. The slope is small, and the cosmological viscosity is large, so the drift is slow rather than rapid.}
\end{classroomqa}

While \(\phi\) inches along the plateau, \(V(\phi)\) changes very little. For that reason the universe inflates almost as if the vacuum energy were constant. The time constant during this stage is of order
\begin{equation}
H^{-1}\sim 10^{-32}\,\mathrm{s}.
\end{equation}
A rough spoken estimate puts the plateau height near \(10^{-14}\) in Planck units.\editorialnote{This is the spoken ``best bet'' estimate at 01:14:38--01:14:47, not a derived value in the lecture.}

How long did inflation last? That is not known, because the width and shallowness of the plateau are not known. What is known is that the universe was multiplied by \(e\) at least about fifty times, and sixty is the number usually quoted.\editorialnote{Editorial restatement of the spoken e-folding discussion at 01:14:05--01:19:27.} These are e-foldings: each e-fold is one multiplication by \(e\).

When the field reaches the edge, the potential energy changes rapidly. The field falls, its potential energy is converted first into kinetic energy and then into heat and particles, and the universe becomes hot. The name ``reheating'' is slightly misleading in this description, because any earlier thermal energy would already have been diluted away during inflation.

The lower minimum is not at zero. The small value left there is what is associated with the present slow exponential expansion.

\begin{classroomqa}
\lecturetimestamp{01:19:48}
\audiencequestion{Could the plateau be exactly flat, with perturbations somehow knocking the field off the edge?}
\lecturerresponse{No. It cannot be completely flat. It has to have some tilt. If it were too flat, the reheated universe would come out too inhomogeneous.}
\end{classroomqa}

``Big Bang'' is used here in two nearby senses. It can mean the reheating event itself. It can also mean the earliest directly visible stage, namely the surface of last scattering. The stable point is that everything actually observed lies far below the inflationary plateau.

\begin{classroomqa}
\lecturetimestamp{01:45:42}
\audiencequestion{Is the reheating itself the Big Bang?}
\lecturerresponse{If you like, yes. But the universe was still extremely hot and opaque at that point. The surface of last scattering comes much later, after the universe has cooled enough to become transparent.}
\end{classroomqa}

\section{Quantum Fluctuations and the Seeds of Galaxies}

Inflation makes the visible universe very homogeneous, but not perfectly homogeneous. With no quantum fluctuations the field would become completely smooth. With quantum mechanics it still varies slightly from place to place.

The central picture plots ordinary space vertically and the field value \(\phi\) horizontally. If the field were exactly homogeneous, every point in space would have the same \(\phi\) at a given time, and every point would drift together toward the edge. Then every region would fall over the cliff at the same moment. Quantum fluctuations spoil that exact lockstep.

Two pictures make the point. One is a small bump in an otherwise homogeneous field. The other is a row of soldiers walking arm in arm toward a cliff. If they are perfectly aligned, they all fall together. If some are slightly ahead and some slightly behind, they go over at different times.

Now comes the important effect. A region that falls over earlier reheats earlier. By the time a later region falls over, the earlier one has already had time to expand and cool somewhat. That difference produces variations in the energy density after reheating.

\begin{classroomqa}
\lecturetimestamp{01:34:00}
\audiencequestion{Are the later reheated regions hotter or cooler than the earlier ones?}
\lecturerresponse{The earlier ones have already had time to expand and cool. The later ones are therefore hotter when they fall over, and that difference contributes to the density inhomogeneity.}
\end{classroomqa}

The observed effect is small:
\begin{equation}
\frac{\delta\rho}{\rho}\sim 10^{-5},
\end{equation}
about one part in one hundred thousand. These are primordial density fluctuations, not dark energy. This number is not derived here from the qualitative sketch alone; it is used to constrain how high and how flat the plateau can be.

They are also necessary. Without them there would be no galaxies. After reheating, gravity amplifies the initial pattern. Over-dense regions attract matter from under-dense regions. The over-dense regions become more over-dense, the under-dense ones more under-dense, and the result is the later structure of galaxies and stars.

This is not only a qualitative story. One can calculate the quantum fluctuations of a field in an expanding universe, follow them through reheating, and compare the result with the fluctuation pattern reconstructed from the galaxy distribution. The fit is very good.

The same discussion also leads to a point about entropy. Gravitational clumping looks counterintuitive, but it does not violate the second law. Gravity is peculiar in statistical mechanics: it tends to make things clump, and that clumping raises the entropy rather than lowering it.

\begin{classroomqa}
\lecturetimestamp{01:40:24}
\audiencequestion{How is the microwave background reflected in all of this, given that \(\phi\) is not an electromagnetic field?}
\lecturerresponse{The primordial density pattern becomes imprinted on the surface of last scattering. Reheating converts the field energy into ordinary particles and photons, and the microwave background later carries the imprint of those primordial fluctuations.}
\end{classroomqa}

\section{Acoustic Waves and the Next Step}

At the end the discussion turns briefly to acoustic waves. These come before the later gravitational clumping. The primordial fluctuations are processed into acoustic oscillations, and those oscillations leave lumps of known physical size on the sky.

That matters because a known physical size at a known distance gives a measuring rod. Once one measures the angle it subtends on the surface of last scattering, geometry can be used to tell whether space is flat, positively curved, or negatively curved. The details are postponed.

\begin{classroomqa}
\lecturetimestamp{01:47:19}
\audiencequestion{What about the acoustic waves?}
\lecturerresponse{They provide a measuring rod on the sky. Their known physical size, together with their observed angular size on the surface of last scattering, is what lets one measure the geometry of space. The details are deferred to the next lecture.}
\end{classroomqa}
